\section{HTCondor Submit-File Generation (\texttt{generators/\allowbreak{}condor/\allowbreak{}})}\label{ch:8}

\subsection{Assembly model}

\texttt{generate\_\allowbreak{}condor\_\allowbreak{}card()} calls one generator per section and joins their
strings in a fixed order: header, retry policy, requirements, undesired sites, authentication, hardware,
executable, file transfer, and queue (always last). Each generator is a pure function of the scard plus
optional overrides, so a submission's card is fully determined by its scard and the call parameters. Adding a
section means adding a function and one line to the assembly list.

\subsection{Header}

\texttt{create\_\allowbreak{}header()} sets \texttt{Universe =\allowbreak{} vanilla} (the only universe
supported for OSG pilots), the CVMFS-hosted Singularity image (\texttt{clas12software:\allowbreak{}production},
or \texttt{:\allowbreak{}devel} when \texttt{devel=\allowbreak{}True}), \texttt{+SingularityBindCVMFS
=\allowbreak{} True} to expose all of \texttt{/\allowbreak{}cvmfs} inside the container, and a \texttt{Rank}
expression built from a site-preference table (\texttt{SITE\_\allowbreak{}RANKS}; e.g. Glasgow 300, CNAF 200,
default 100, Syracuse deprioritized).

\subsection{Retry / hold policy}

\texttt{create\_\allowbreak{}retry\_\allowbreak{}policy()} encodes eviction tolerance:
\texttt{on\_\allowbreak{}exit\_\allowbreak{}remove} removes only clean exits (no signal, code 0);
\texttt{on\_\allowbreak{}exit\_\allowbreak{}hold} holds any signalled or non-zero exit;
\texttt{periodic\_\allowbreak{}release} releases a held job back to idle only when it has been retried fewer
than 5 times, has been held over an hour, and exited with one of the known transient OSG codes (202, 204--208,
210--213 --- CVMFS unavailable, pilot evicted, etc.). Anything else stays held for inspection. The exit-code
taxonomy is in Appendix E.

\subsection{Requirements}

\texttt{create\_\allowbreak{}requirements()} guards against known infrastructure failures:
\texttt{HAS\_\allowbreak{}SINGULARITY}, the OSG oasis CVMFS repository mounted, kernel $\geq$ 2.17, a
sufficiently recent CVMFS oasis revision, and glidein version $\geq$ 534.
\texttt{-\allowbreak{}-\allowbreak{}target-\allowbreak{}site} appends a \texttt{GLIDEIN\_\allowbreak{}Site
=\allowbreak{}=\allowbreak{} "<SITE>"} clause to pin all subjobs to one site.

\subsection{Undesired sites}

\texttt{create\_\allowbreak{}undesired()} emits a \texttt{+UNDESIRED\_\allowbreak{}Sites} exclusion list,
keeping jobs off sites known to fail for CLAS12 workloads.

\subsection{Authentication}

\texttt{create\_\allowbreak{}authentication()} emits \texttt{use\_\allowbreak{}oauth\_\allowbreak{}services
=\allowbreak{} jlab\_\allowbreak{}clas12}. The submit node's CredMon obtains an OAuth2 bearer token for that
service; HTCondor injects it as \texttt{BEARER\_\allowbreak{}TOKEN\_\allowbreak{}FILE} into the job
environment, where Pelican/OSDF use it to read LUND input and write HIPO output. No long-lived secret ever
reaches the node.

\subsection{Hardware}

\texttt{create\_\allowbreak{}hardware()} requests the GEMC resource profile: 1 CPU (GEMC is single-threaded),
2.256 GB memory (GEMC peak RSS $\approx$ 1.8 GB plus headroom), and 2 GB scratch disk. Overrides exist for
heavier workloads.

\subsection{Executable}

\texttt{create\_\allowbreak{}executable()} names \texttt{nodescript.\allowbreak{}sh} as the executable, sets
the per-subjob \texttt{log}/\texttt{out}/\texttt{err} paths under \texttt{log/\allowbreak{}}, and sets
\texttt{+ProjectName} for OSG accounting.

\subsection{File transfer}

\texttt{create\_\allowbreak{}file\_\allowbreak{}transfer()} stages \texttt{nodescript.\allowbreak{}sh} and
\texttt{functions.\allowbreak{}sh} in (\texttt{should\_\allowbreak{}transfer\_\allowbreak{}files =\allowbreak{}
YES}, \texttt{when\_\allowbreak{}to\_\allowbreak{}transfer\_\allowbreak{}output =\allowbreak{}
ON\_\allowbreak{}EXIT}) and retrieves only the \texttt{output/\allowbreak{}} subdirectory. Notably, type-2 LUND
files are \textbf{not} staged through HTCondor: the worker node fetches each file on-node via Pelican, avoiding
saturation of the submit node with large LUND directories.

\subsection{Queue}

\texttt{create\_\allowbreak{}queue()} is the final section. Type-1 emits \texttt{Arguments =\allowbreak{}
\$(Process)} and \texttt{Queue N} (N = \texttt{njobs}). Type-2 uses itemdata: \texttt{Arguments =\allowbreak{}
\$(Process) \$(lundFile)} and \texttt{queue lundFile from lund\_\allowbreak{}files}, expanding the staged
\texttt{lund\_\allowbreak{}files} one line per subjob.

\subsection{Extending the card}

Because each section is an independent pure function joined in order, a new section (say, a new accounting
attribute or a new requirement) is added by writing one generator and inserting one call in
\texttt{generate\_\allowbreak{}condor\_\allowbreak{}card()} --- no template surgery. An annotated full card is
in Appendix C.
